\begin{solapartat}
\punts{0,5} Segons la llei de Coulomb, el mòdul de la força sobre l'electró degut a l'atracció del protó és:
\[ F_e = k\,\frac{e^2}{r^2} \]
La segona llei de Newton estableix que $\vec{F} = m\vec{a}$.

D'altra banda, considerant que la trajectòria és circular, la seva acceleració és l'acceleració centrípeta, $a = v^2/r$:
\[ k\,\frac{e^2}{r^2} = \frac{m\cdot v^2}{r} \Rightarrow m\cdot v^2 = k\,\frac{e^2}{r} \]
De l'anterior expressió, podem obtenir-ne la velocitat:
\[ v = e\sqrt{\frac{k}{m\cdot r}} \]
\punts{0,25} El moment angular s'expressa així:
\[ \vec{L} = m\vec{r}\times\vec{v} \]
En què $\vec{r}$ és el vector de posició que apunta del protó a l'electró. El mòdul del moment angular s'expressa així:
\[ L = m\,r\,v\,\sin\theta \]
En què $\theta$ és l'angle que formen $\vec{r}$ i $\vec{v}$. Com que l'òrbita és circular, $\vec{r}$ i $\vec{v}$ són perependiculars, $\theta = \pi/2$ i $\sin\theta = 1$.

Per tant:
\[ L = m\cdot r\cdot v = m\cdot r\cdot e\sqrt{\frac{k}{m\cdot r}} = e\sqrt{k\cdot m\cdot r} \]
\punts{0,25} Si fem servir l'expressió $L = \frac{n\cdot h}{2\pi}$:
\[ \frac{n\cdot h}{2\pi} = e\sqrt{k\cdot m\cdot r} \Rightarrow r = \frac{1}{k\cdot m}\left(\frac{n\cdot h}{2\pi\cdot e}\right)^2 \]
Alternativament, es pot plantejar $L = m\cdot r\cdot v = \frac{n\cdot h}{2\pi}$ i $v = e\sqrt{\frac{k}{m\cdot r}}$.

\punts{0,25} I el radi per a $n = 1$ és:
\begin{align*}
r = \frac{1}{k\cdot m}\left(\frac{n\cdot h}{2\pi\cdot e}\right)^2
&= \frac{1}{8,99\cdot 10^{9}\cdot 9,11\cdot 10^{-31}}\left(\frac{6,63\cdot 10^{-34}}{2\pi\cdot 1,602\cdot 10^{-19}}\right)^2\\
&= 5,30\cdot 10^{-11}\ \mathrm{m}
\end{align*}
\end{solapartat}

\begin{solapartat}
De l'apartat anterior, tenim:
\[ v = e\sqrt{\frac{k}{m\cdot r}} = 2,18\cdot 10^{6}\ \mathrm{m/s} \]
\destaca{Alternativament:} $L = e\sqrt{k\cdot m\cdot r} = 1,055\cdot 10^{-34}\ \mathrm{kg}\,\frac{\mathrm{m^2}}{\mathrm{s}} \Rightarrow v = \frac{L}{m\cdot r} = 2,18\cdot 10^{6}\ \mathrm{m/s}$ o

$L = \frac{h}{2\pi} = 1,055\cdot 10^{-34}\ \mathrm{kg}\,\frac{\mathrm{m^2}}{\mathrm{s}} \Rightarrow v = \frac{L}{m\cdot r} = 2,18\cdot 10^{6}\ \mathrm{m/s}$

\punts{0,5} I l'energia cinètica és:
\[ E_c = \frac{1}{2}m\cdot v^2 = \frac{1}{2}9,11\cdot 10^{-31}\cdot (2,18\cdot 10^{6})^2 = 2,18\cdot 10^{-18}\ \mathrm{J} \]

\punts{0,5} L'energia potencial s'expressa així:
\[ E_p = -k\,\frac{e^2}{r} = -8,99\cdot 10^{9}\,\frac{(1,602\cdot 10^{-19})^2}{5,3\cdot 10^{-11}} = -4,35\cdot 10^{-18}\ \mathrm{J} \]

\punts{0,25} I l'energia mecànica:
\[ E_m = E_c + E_p = 2,18\cdot 10^{-18} - 4,35\cdot 10^{-18} = -2,18\cdot 10^{-18}\ \mathrm{J} \]
\end{solapartat}
